\section{Discussion}
\label{sec:discussion}

\paragraph{Why stale coherence dissociates from existing memory
benchmarks.}
LongMemEval~\citep{Wu2025LongMemEval} and RULER~\citep{Hsieh2024RULER}
ask ``did the model retrieve the fact?'' --- they probe only the
verbal axis. MemoryArena~\citep{He2026MemoryArena} and
MemoryAgentBench~\citep{Hu2025MemoryAgentBench} probe action
grounding, but they do so inside session-discrete loops in which
the agent receives explicit reset boundaries between trials. Stale
coherence is visible only when both axes are measured \emph{jointly}
on tick-continuous, naturally-accumulated context --- i.e., on the
trajectory shape that production agent deployments actually exhibit,
in which an anchor injected at deployment time has to survive hours
of intermediate decisions without being re-prompted. This is
exactly the regime our 8-sim-hour episode and dual probe
($\texttt{verbal\_recall} \times \texttt{action\_grounded\_recall}$)
target.

\paragraph{Implications for production agent design.}
Three concrete recommendations follow:
\textbf{(i)} Test action-grounded recall, not summary fidelity. A
model that scores well on LongMemEval may still abandon anchors in
its action stream within 4 sim-hours; if your deployment cares
whether the agent \emph{acts on} its memories rather than merely
mentions them, verbal-only benchmarks will systematically
over-estimate readiness.
\textbf{(ii)} The verbal-recall axis is necessary but not
sufficient for memory benchmark coverage. Existing benchmarks
should add an action-grounded probe wherever the underlying task
exposes structured directives or tool calls, so that
$\Delta(t)$ can be reported alongside the verbal score.
\textbf{(iii)} If your agent must maintain anchor commitment across
hours of deployment, monitor $\Delta(t)$ in production telemetry,
not just summary mentions. A flat verbal-recall curve combined
with a decaying action-grounded-recall curve is the operational
signature of stale coherence in the wild.

\paragraph{What this implies for ``memory'' research.}
Stale coherence suggests that LLM ``memory'' is not one capability
but at least two --- verbal echo and behavioural commitment ---
that can decouple under naturally-accumulated context. Future memory
benchmarks should report both axes; the
\citet{Sinha2025IllusionDiminishingReturns} curve relating model
scale to long-context recall may need to be reported separately
for each, because a regime in which verbal recall improves with
scale while action-grounded recall does not would be invisible to
single-axis evaluation but visible to the dual probe defined here.
